\documentclass[a4paper,11pt]{ltjsarticle}
\usepackage[haranoaji]{luatexja-preset}
\usepackage{amsmath,amssymb}
\usepackage[top=12mm,bottom=10mm,left=14mm,right=14mm]{geometry}
\usepackage{array}
\usepackage{tikz}
\usetikzlibrary{calc}
\pagestyle{empty}
\setlength{\parindent}{0pt}
\newlength{\qcol}\setlength{\qcol}{0.47\textwidth}
\newlength{\qgap}
\newlength{\anscolw}\setlength{\anscolw}{0.30\textwidth}
\newlength{\ansh}
\newcommand{\Q}[2]{\parbox[t]{\qcol}{\raggedright (#1)\hspace{0.6em}$\displaystyle #2$}}
\newcommand{\QR}[4]{\Q{#1}{#2}\hfill\Q{#3}{#4}\par\vspace{\qgap}}
\newcommand{\anscell}[1]{\rule[-\ansdrop]{0pt}{\ansh}(#1)}
\newlength{\ansdrop}

\begin{document}
{\large\bfseries 計算問題　19}\hfill 名前(\underline{\hspace{32mm}})\par\vspace{3mm}
■（1）〜（16）の計算をしなさい。（17）〜（21）は方程式を解きなさい。\par\vspace{3mm}
\setlength{\qgap}{5.4mm}
\QR{1}{3(2x-5)-(2x-7)}{2}{(2a+3)+(3a-1)}
\QR{3}{28x\div(-4)}{4}{(-11)\times13a}
\QR{5}{8-(-5)+(-3)^2}{6}{-6+(-19)}
\QR{7}{\dfrac{x-1}{4}\times8}{8}{(24x-8)\div\dfrac{2}{3}}
\QR{9}{(-3)^2\times(-2)^3}{10}{(-6)-(-11)}
\QR{11}{2-4+6-8}{12}{3\div7\times(-14)}
\QR{13}{-3^2+12\div(-4)}{14}{\dfrac{3}{4}(4x-12)+\dfrac{2}{3}(6x-3)}
\QR{15}{(-5)\times3^2+(-2)^3}{16}{-24\div\{-3\times(5-7)\}}
\QR{17}{-2(x+7)=-3(x+4)}{18}{-7x-8=13}
\QR{19}{0.1x+0.6=2}{20}{\dfrac{5}{3}x+2=\dfrac{3}{2}x}
\vspace{1mm}
\Q{21}{0.4(x-3)=0.2x+1}\par\vspace{3mm}
\setlength{\ansh}{9.5mm}\setlength{\ansdrop}{6mm}%
\begin{center}\begin{tabular}{|p{\anscolw}|p{\anscolw}|p{\anscolw}|}\hline
\anscell{1} & \anscell{2} & \anscell{3} \\\hline
\anscell{4} & \anscell{5} & \anscell{6} \\\hline
\anscell{7} & \anscell{8} & \anscell{9} \\\hline
\anscell{10} & \anscell{11} & \anscell{12} \\\hline
\anscell{13} & \anscell{14} & \anscell{15} \\\hline
\anscell{16} & \anscell{17} & \anscell{18} \\\hline
\anscell{19} & \anscell{20} & \anscell{21} \\\hline
\end{tabular}\end{center}

\end{document}
